\section{Mathematical Overview}
\seclabel{new-overview}
\emph{Universal algebra} concerns the generic description of algebraic
structures and relationships between them~\cite[Chp.~II]{burris-sankappanavar:universal-algebra}.
We summarise
the concepts relevant to \Frexlib{}.

\subsection{Presentations of Algebraic Structures}\seclabel{presentations}
A \emph{finitary signature} $\Sig = \pair{\Op\Sig}{\arity}$ consists
of a set $\Op\Sig$ of \emph{operators},
and an assignment $\arity : \Op\Sig \to \Nat$ of a natural
number to each operator called its \emph{arity}.
%
For example, the signature $\Sig_{(2,0)}$ often used for monoids
has two operators $\Op(\Sig_{(2,0)}) \definedby \set{(+),0}$ with
respective arities $2$ and $0$.
%
It is common to use a succinct notation that groups operators
and their arities: $\Sig_{(2,0)} \definedby \set{(+):2,0:0}$.

Signatures determine an algebraic language, whose semantic models
are \emph{algebras}.
%
An \emph{algebra} $\A = \pair{\U\,\A}{\A\sem-}$
  for a signature $\Sig$
  consists of
a set $\U\,\A$ called the \emph{carrier} and
an assignment $\A\sem f : (\U\,\A)^n \to \U\,\A$ of an \emph{$n$-ary operation} over this carrier for every $n$-ary operator $f : n$ in $\Sigma$.
%
Continuing the example above, $\Sig_{(2,0)}$-algebras amount to
triples $\triple{X}{\sem{(+)} : X^2 \to X}{\sem{0} \in X}$.
%
There may be many different algebras for a given signature and carrier set:
for instance, we can equip the natural
numbers $\naturals$ with the $\Sig_{(2,0)}$-algebra structures of
arithmetic addition $\triple{\naturals}{(+)}{0}$,
arithmetic multiplication $\triple\naturals{(\cdot)}1$,
or $\triple{\naturals}{\max}{0}$.
%
Examples abound:
subtraction over the integers has a $\Sig_{(2,0)}$-algebra
$\triple{\integers}{(-)}{0}$;
$n\times n$ matrices over $\naturals$ have
$\Sig_{(2,0)}$-algebra structures given by matrix addition and
multiplication with the zero and identity matrix respectively, and so
on.
%

Each signature determines a language consisting of \emph{terms}.
Given a set $\X$ of variables, the
\emph{$\Sig$-terms over $\X$} are given inductively as either a
variable in $\X$ or an application $f(t_1, \ldots, t_n)$ of an
operator $f : n$ from $\Sig$ to $n$ terms over $\X$.
The primary role of terms is to designate \emph{equations in context}
$\X \vdash t = s$, i.e.~triples consisting of a set $\X$ of
variables and two terms in context $\X$. For example, the
associativity equation, expressed over the signature $\Sig_{(2,0)}$,
is $x,y,z \vdash x + (y + z) = (x + y) + z$.

An \emph{environment} for a context (=set) $\X$ in an algebra $\A$ is
a function $\env : \X \to \U\,\A$.
%
An algebra $\A$ determines, for
each term in context $\X \vdash t$, an \emph{interpretation} function
$\A\sem t : (\U\,\A)^{\X} \to \U\,\A$.
%
Given an environment $\env$, $\A\sem t$ interprets
each variable in the environment $\A\sem{\x}\env \definedby \env\,\x$,
and structurally interprets each operator
application:
$\A\sem{f(t_1, \ldots, t_n)}\env\definedby \A\sem f(\A\sem{t_1}\env, \ldots, \A\sem{t_n}\env)$.
For example, the interpretation of the left-hand-side (LHS) of the
associativity axiom in the $\Sig_{(2,0)}$-algebra $\triple{\naturals}{\max}{0}$
given the environment $\{x \mapsto 5, y \mapsto 3, z \mapsto 8\}$ is
$\max(5, \max(3, 8)) = 8$.

We say that an equation is \emph{valid} in an algebra $\A$, writing
$\A \models (\X \vdash t = s)$, when $\A\sem t\env = \A\sem s\env$
for all environments $\env : \X \to \U\,\A$. %% It is this implicit universal
%% quantification over the environment that gives universal algebra its name.
For example, the $\Sig_{(2,0)}$-algebras
presented so far validate the associativity axiom, but interpreting
the binary operation as subtraction over the integers $\triple{\integers}{(-)}{0}$
does not: taking the
environment $\{x \mapsto 0, y \mapsto 0, z \mapsto 1\}$, we have
$0 - (0 - 1) = 1 \neq -1 = (0 - 0) - 1$.

A \emph{presentation} of an algebraic theory $\Pres = \pair{\Sig_{\Pres}}{\Pres\Axiom}$
consists of a signature $\Sig_{\Pres}$ and a set $\Pres\Axiom$ of
$\Sig_{\Pres}$-equations in context, which we call \emph{axioms}. A
\emph{$\Pres$-model} $\A$ is a $\Sig_{\Pres}$-algebra $\A$ validating
all $\Pres$-axioms.  For example, the axioms of the $\Monoid$
presentation consist of associativity and neutrality ($x \vdash x
+ 0 = x, 0 + x = x$) over the signature $\Sig_{(2,0)}$. The
axioms for the $\CMonoid$ presentation additionally include commutativity $x,y
\vdash y + x = x + y$. We can now generically manipulate classes of
algebraic structures using these concepts, while generalising the
usual examples: $\Monoid$-models are monoids, $\CMonoid$-models are
commutative monoids, etc.

\subsection{Homomorphisms, Free Models/Algebras, and Free Extensions}\seclabel{simplification}
For a presentation $\Pres$, we consider two classes of simplification problem of interest, which we call \emph{purely abstract} and \emph{partially concrete}, respectively.
%
In the purely abstract case, the input to the problem consists of a set of variables $\X$ and a $\Sigma_{\Pres}$-term over $\X$.
%
Usually, $\X = \Finn$ is a set of De Bruijn levels.
%
In the partially concrete case, the input instead consists of a $\Pres$-model $\A$, a set of variables $\X$ and a $\Sigma_{\Pres}$-term over the disjoint union $\U\,\A \uplus \X$.
%
Here, we do not treat the `variables' in the term coming from the inclusion of $\U\,\A$ as variables in the usual sense.
%
Indeed, in both cases, the goal of the simplification problem is to find a representative modulo the presentation's axioms and the rules of deduction.
%
However, in the partially concrete case, simplification must also fold the elements of $\U\,\A$ using the operations of $\A$.
%
When the set $\X$ is empty, simplification and evaluation should coincide.

The foregoing description is imprecise, but can be made precise using the notions of free algebra and free extension, whose definitions are formulated in terms of the unique existence of certain structure-preserving maps.
%
We dedicate the remainder of this section to introducing these concepts, drawing connection to the frex and fral simplification examples from
page \pageref{new-intro fral example}. We start with the fral:
\begin{equation}
-6 + (x+3) + (y+x)
{}\xmapsto{\text{evaluate}(x_0 + (x_2 + x_1) + (x_3 + x_2))}{}
(1, 1, 2, 1)
{}\xmapsto{\text{reify}\left(\begin{array}{@{}>{\scriptstyle}l@{}}
      x_0 \mapsto -6, x_1 \mapsto 3, \\
      x_2 \mapsto x, x_3 \mapsto y
    \end{array}
\right)}{}
-6 +3 + 2x + y
\qquad
\label{nbe-examples-fral}
\end{equation}

Let $\Pres$ be a presentation and $\X$ a set of variables.  We define a \emph{$\Pres$-model\/}
\emph{$\avar{}=\pair{\aModel}{\aEnv}$ over $\X$} to be a
$\Pres$-model $\aModel$ equipped with an $\X$-environment in this
algebra, i.e.~a function $\env : \X \to \U(\aModel)$.  This concept
formalises the inputs to the fral simplification process from
\eqref{nbe-examples-fral}.
The starting term is the expression $-6 + (x+3) + (y+x)$
in the meta-language involving the
meta-level variables $x$ and $y$.
The argument
$x_0 + (x_2 + x_1) + (x_3 + x_2)$
in the label on the left arrow is
a $\Sig_{(2,0)}$-term with object-level variables from
$\X \definedby \set{x_0, \ldots, x_3}$. The argument to the label on
the right arrow, $(x_0 \mapsto -6, x_1 \mapsto 3, x_2 \mapsto x, x_3
\mapsto y)$, is an $\X$-environment mapping the object-level variables $x_i \in \X$
to expressions in the meta-level such as the constants $-6$ and $3$ and the
meta-level variable $y$. To the fral simplification process
these three expressions have the same status.
The frex simplification process we describe later uses the concrete
nature of $-6$ and $3$, simplifying the meta-level expression further through evaluation.

Let $\A[1]$, $\A[2]$ be $\Sig$-algebras for a signature
$\Sig$.  A \emph{homomorphism} $h : \A[1] \to \A[2]$ of
$\Sig$-algebras is a semantics-preserving function $h : \U\,\A[1]
\to \U\,\A[2]$ between their carriers. Explicitly, for all operators
$f : n$ in $\Sig$ and elements $a_1, \ldots, a_n$ in $\U\,\A$, we have:
\( h(\A[1]\sem f\!(a_1, \ldots, a_n)) = \A[2]\sem f\!(h\,a_1,
\ldots, h\,a_n) \). A homomorphism between $\Pres$-models
is a homomorphism between the underlying signature algebras.
For example, the list-length function is a homomorphism from the
monoid of concatenation over lists to the monoid
of addition over naturals: $\length : \triple{\List \X}{(\++)}{\Empty}
\to \triple\naturals{(+)}0$.

\begin{wrapfigure}[4]{r}{2cm}
  \vspace{-1\baselineskip}
  \includegraphics{diagram-01.mps}
\end{wrapfigure}
A \emph{morphism} $\hName : \avar \to \bvar$ of $\Pres$-models over
$\X$ is a $\Pres$-model homomorphism that moreover makes the diagram
on the right commute.  A \emph{free $\Pres$-model over $\X$} is
then a $\Pres$-model over $\X$ from which there is a unique such
morphism to every other $\Pres$-model over $\X$. This
existence-and-uniqueness property is called the \emph{universal
property} of free $\Pres$-models.  For example, the free commutative
monoid over $\X = \set{x_0, \ldots, x_3}$ is the $\Sig_{(2,0)}$-algebra
over $\naturals^4$ given by componentwise arithmetic addition, and
equipped with the $\X$-environment sending $x_0$ to $(1, 0, 0, 0)$,
etc. The unique morphism out of this algebra into any commutative
monoid $\A$, equipped with an environment $\env$, sends $(a_0, \ldots,
a_3)$ to $a_0(\env\,x_0) + \ldots + a_3(\env\,x_3)$.
%

Confusingly, the literature sometimes refers to $\Pres$-models as
$\Pres$-algebras~\cite[\S{}V.6]{mac-lane:categories-2nd-ed}. To avoid
this confusion, we will use the terminology `$\Pres$-models' as much as
possible, with one exception. When $\Pres$ is implicit, we will
sometimes use the portmanteau `fral' of `\textbf{fr}ee
\textbf{al}gebra'. We prefer it over alternatives (`frmo' or `frem').

In summary, fral simplification provides a guiding principle
for designing the data structure in the middle of
\eqref{nbe-examples-fral}: it is an implementation of the fral.
%
\label{methodology-anchor}%
The success criterion for the design of the simplification code is the
universal property, which singles the fral up to a unique isomorphism
of models over $\X$. The different components of the universal
property provide the following methodology for this creative process.  Failure in
each step in this process may provide insight into how the current
data structure fails, and how it may be revised:
\begin{enumerate}
\item\label{methodology-step-structure} \textit{Equip the data structure with
  a $\Sig_{\Pres}$-algebra structure.} Failure in this step typically exposes
  inputs this simplifier cannot deal with. Revise the data structure to
  support these operations.
\item\label{methodology-step-validation} \textit{Prove it validates the
  axioms, i.e., forms a $\Pres$-model.}  Failure in this step
  exposes equations the structure does not use, i.e., completeness
  issues. Revise the data structure to use the additional equations, or
  relax the presentation, giving up on using these equations.
\item\label{methodology-step-existence} \textit{Define a function out of
  it to any other $\Pres$-model over $\X$.} Failure in this step
  typically exposes `junk' in the data structure: components that are
  not captured by the presentation such as missing operators.
  Revise the data structure to exclude the junk, or extend the presentation
  with more operators to account for the additional components.
\item\label{methodology-step-homomorphic}
  \textit{Show this function is a homomorphism.}
  Failure in this step typically exposes unprovable
  equations that the data structure validates.
  Revise the data structure/presentation.
\item\label{methodology-step-4-uniqueness} Show this homomorphism is
  unique.  Failure in this step similarly exposes implicit axioms in
  the data structure, revise data structure/presentation.
\end{enumerate}
While each step requires mathematical creativity, the overall structure
breaks the task down into smaller steps, a checklist. This checklist
organises the simplification code along these steps, making simplifier
development more methodical.

We now turn to frex simplification, through the second example from
page~\pageref{new-intro fral example}:
\begin{equation}
-6 + (x+3) + (y+x)
{}\xmapsto{\text{evaluate}_{\integers}}{}
(-3, 2, 1)
{}\xmapsto{\text{reify}}{}
-3 + 2x + y
\label{nbe-examples-frex}
\end{equation}
Let $\A$ be a $\Pres$-model and $\X$ a set whose elements we treat as representing
variables. An \emph{extension}
of $\A$ by $\X$ is a triple
$\avar =
\triple{\aModel}{\aVar}{\aEmbed}$
consisting of a $\Pres$-model $\aModel$, a function $\aVar : \X \to
\U(\aModel)$, and a $\Pres$-homomorphism $\aembed : \A \to \aModel$.
This concept lets us phrase the distinction between meta-level constants such as
$3$ and $6$, and meta-level variables such as $y$ in the meta-level expression
$-6 + (x+3) + (y+x)$ from \eqref{nbe-examples-frex}.
Taking $\A := \integers$ and $\X \definedby \set{x,y}$, both the
concrete elements of the algebra $\A$ (e.g., $3$ and $-6$) and the
abstract variables from $\X$ are part of the vocabularly of the
simplification process. It uses the evaluation equation
$-6 + 3 = -3$.

\begin{wrapfigure}[4]{r}{3.2cm}
  %uncomment next line if this paragraph is not in the beginning of the page.
  \vspace{-1\baselineskip}
  \includegraphics{diagram-02.mps}
\end{wrapfigure}
A \emph{morphism}\/ $\hName : \avar \to \bvar$ of extensions of $\A$
by $\X$ is a $\Pres$-homomorphisms $\hName : \aModel \to \bModel$ that
makes the two triangles in the diagram on the right commute.  A
\emph{free extension} of $\A$ by $\X$ is then an extension from which
there is a unique such morphism to every other extension of $\A$ by
$\X$.  For example, the free extension of a commutative monoid $\A$ by
two variables $\X \definedby \set{x,y}$ is the $\Sig_{(2,0)}$-algebra
over $\U\,\A\times \naturals^2$, given componentwise, equipped with
the $\Sig_{(2,0)}$-homomorphism that sends $u \in \U\,\A$ to $(u, 0,
0)$ and the function that sends $x$ to $(\A\sem0, 1, 0)$ and $y$ to
$(\A\sem0, 0, 1)$. The unique morphism to any
extension $(\A[2], \env, b)$ sends $(u,a_0,a_1)$ to
$b u + a_0(\env\,x) + a_1(\env\,y)$.

\begin{figure*}
  \centering
  \begingroup\small
\begin{tikzpicture}

\node[minimum width=0.3\textwidth] (monoid)
  at (0,0) { \bfseries monoids };
\node[minimum width=0.3\textwidth,anchor=north west] (cmonoid)
  at (monoid.north east) { \bfseries \textit{commutative} monoids };
\node[minimum width=0.3\textwidth,anchor=north west] (imonoid)
  at (cmonoid.north east) { \bfseries \textit{involutive} monoids };

\node[anchor=north] 
  at ($(monoid.north) + (0,-0.8)$)
     { \parbox{0.28\textwidth}{\centering lists of\\ variables} };
\node[anchor=north]
  at ($(cmonoid.north) + (0,-0.8)$)
     { \parbox{0.28\textwidth}{\centering origin-intercepting\\ linear polynomials} };
\node[anchor=north]
  at ($(imonoid.north) + (0,-0.8)$)
     { \parbox{0.28\textwidth}{\centering lists over ordinary \& \\ singly-involuted variables } };

\node[anchor=north]
  at ($(monoid.north) + (0,-1.8)$)
     { $yxxyx$ };
\node[anchor=north]
  at ($(cmonoid.north) + (0,-1.8)$)
     { \parbox{0.28\textwidth}{\centering $a_1x_1\! +\! \ldots\! +\! a_nx_n$ \\ $(a_i : \Nat)$}};
\node[anchor=north] 
  at ($(imonoid.north) + (0,-1.8)$)
     { $y\overline xxx\overline yx$ };

\node[anchor=north]
  at ($(monoid.north) + (0,-2.9)$)
     { \parbox{0.28\textwidth}{\centering ~\\[-0.5ex] alternating lists\\in $\mathbb{M}_{2\times 2}(\Nat{})[y]$\\~} };
\node[anchor=north]
  at ($(cmonoid.north) + (0,-2.9)$)
     { \parbox{0.28\textwidth}{\centering ~\\[-0.5ex]  linear polynomials \\ in $\A{}[x_1, \ldots, x_n]$\\~} };
\node[anchor=north]
  at ($(imonoid.north) + (0,-2.9)$)
     { \parbox{0.28\textwidth}{\centering alternating lists\\ with tagged variables \\ in $\String{}[x,y]$ } };

\node[anchor=north]
  at ($(monoid.north) + (0,-4.1)$)
     {   $\begin{pmatrix} 1 & 3 \\ 0 & 2 \end{pmatrix}y\begin{pmatrix} 0 & 1 \\ 1 & 0 \end{pmatrix}y$ };
\node[anchor=north]
  at ($(cmonoid.north) + (0,-4.3)$)
     { \parbox{0.28\textwidth}{\centering $c\! +\! a_1x_1\! +\! \ldots\! +\! a_nx_n$ \\ ($a_i : \Nat{}, c : \A$)}};
\node[anchor=north] (imonoid-frex-repr-eg)
  at ($(imonoid.north) + (0,-4.3)$)
     { $\EmptyStr{}x\helloStr{}y\ollehStr{}\overline x\EmptyStr{}$ };

\node[rectangle,rounded corners,draw=black,minimum width=0.26\textwidth,minimum height=5.3cm,anchor=north west]
  at ($(monoid.north west) + (0.01\textwidth,0)$) { };
\node[rectangle,rounded corners,draw=black,minimum width=0.26\textwidth,minimum height=5.3cm,anchor=north west]
  at ($(cmonoid.north west) + (0.01\textwidth,0)$) { };
\node[rectangle,rounded corners,draw=black,minimum width=0.26\textwidth,minimum height=5.3cm,anchor=north west]
  at ($(imonoid.north west) + (0.01\textwidth,0)$) { };

\node[rectangle,fill=red,fill opacity=0.1,minimum width=0.96\textwidth,minimum height=2cm,anchor=north west] (frals)
  at ($(monoid.north west) + (-0.06\textwidth,-0.7cm)$) { };

\node[rectangle,fill=green,fill opacity=0.1,minimum width=0.96\textwidth,minimum height=2.2cm,anchor=north west] (frexes)
  at ($(monoid.north west) + (-0.06\textwidth,-2.9cm)$) { };

\node[anchor=north,rotate=90] at ($(frals.west) + (0.03\textwidth,0)$) { free algebra};
\node[anchor=north,rotate=90] at ($(frexes.west) + (0.03\textwidth,0)$) { free extension};

\end{tikzpicture}
\endgroup

  \caption{Frals and frexes for varieties of monoids}
  \figlabel{monoid frexlets}
\end{figure*}
\Figref{monoid frexlets} summarises the frals and frexes
in this article. All but the last frex are well-known representations.
Our contribution is to implement, alongside these representations,
constructive proofs that they represent the fral or the frex. These proofs
require the universal algebraic concepts we introduced
so far. Other simplification modules can reuse these concepts and the representation
theorems for existing simplifiers. This design enables extensible and modular
user and library code.

\subsection{Setoids}

Designing \Frexlib{} around the concepts from
\secref{presentations} and~\secref{simplification} supports
term simplification in a uniform way.
%
Generalising the design further to use
\emph{setoids}~\cite{hofmann:thesis,Bishop1967-BISFOC-3} instead of
sets enables the same abstractions to offer more flexible
functionality: printing terms and proofs, proof simplication, code
generation, and more, which
\secsref{certification}{reflection}
explore.  The ability to extract such
intensional information from terms violates algebraic equality, and
differentiates this design from work involving quotients, for example
quotient inductive-inductive types
(QIITs)~\cite{altenkirch-et-al:qiits,kaposi-et-al:constructing-qiits,altenkirch-et-al:tt-in-tt-with-qiits}.
This
section presents some mathematical background about setoids and their
realisation in \idris{}.

A \emph{setoid} $\X = (\U\,\X,(\equiv_{\X}))$ consists of a set
$\U\,\X$ and an equivalence relation $(\equiv_{\X})$ over $\U\,\X$.  A
\emph{setoid homomorphism} $\f : \X[1] \~> \X[2]$ is a
relation-preserving function between sets $\X[1]$ and $\X[2]$:
$\f\,x \equiv_{\X[2]} \f\,y$ whenever $x \equiv_{\X} y$.
We think of elements in $\U\,\X$ as representatives of the equivalence
classes of $(\equiv_{\X})$, and so every setoid homomorphism induces
a (unique) function between the quotients $\X/(\equiv_{\X}) \to \X[2]/(\equiv_{\X[2]})$.
One way to construct a setoid is to equip a set $\X$ with its
equality relation to give $(\X, (=))$, but using setoids
also allows us to explicate sophisticated equivalence relations and define
operations on them that are not supported by the corresponding
quotients. For example, given a presentation $\Pres$ and a set $\X$,
we can define the \emph{provability} relation $\X \vdash - = -$ over
terms $\X \vdash t$ (cf.~\figref{provability}).  Related elements in
the setoid $(\Term\X, \X\vdash -=-)$ represent different, but provably
equal, terms. In contrast, elements in the quotient
$\Term\X/(\X\vdash-=-)$ represent equivalence classes of provably equal
terms. \Figref{provability} also includes the evaluation axiom (\textsc{eval}),
which we will use to present the frex.
\begin{figure}
  \begin{gather*}
\inferrule
    {~}
    {
      \X \vdash t = t
    }
    \textsc{refl}
\qquad
\inferrule
    {\X \vdash s = t}
    {
      \X \vdash t = s
    }\textsc{sym}
\qquad
\inferrule
    {\X \vdash t = s\quad
      \X \vdash s = r
    }
    {
      \X \vdash t = r
    }\textsc{trans}
\qquad
\inferrule{
  (\X \vdash t = s) \in \Pres\Axiom
}{
  \X \vdash t = s
}\textsc{ax}
\\
 \inferrule
    {\X[2] \vdash t=s\quad
      \theta_1,\theta_2 : \X[2] \to \Term\X[2]\quad
     (\X \vdash \theta_1y = \theta_2y)_{y \in \X[2]}
    }
    {
      \X \vdash t[\theta_1] = s[\theta_2]
    }\textsc{cong}
\qquad
\shade{
\inferrule{
  \U\A \vdash t
}{
  \X \vdash t = \A\sem{t}
}\textsc{eval}
}
\end{gather*}

  \caption{Provability in (a) equational logic (unshaded) and (b) the evaluation rule}
  \figlabel{provability}
\end{figure}

Setoids and their homomorphisms form a common technique to complete an
intensional type theory with extensional functions and quotients,
requiring users to establish that every defined function is a setoid
homomorphism.  However, in \Frexlib{} we make essential use of setoids
that quotient types~\cite{hofmann-phd} do not allow.
%
For example, the terms-up-to-provability
setoid supports operations such as $\mathrm{vars} : \Term\X
\to \List\;\X$ which extract the list of variables appearing in a given
term in-order. This function is not a setoid homomorphism.
The quotient can only support such extraction following a canonicalisation.
Both flavours of
function are useful in applications, but setoids, and not quotients,
support both.

In \Frexlib{}, we implement universal algebra \emph{internal} to setoids:
carriers are setoids; algebraic operations are setoid homomorphisms;
algebra homomorphisms and environments must also be setoid
homomorphisms; and the unique homomorphisms in the universal
properties of the fral and the frex must be setoid homomorphisms.
With this generalisation, the setoid $\Term\X/(\X\vdash-=-)$ also
satisfies the fral universal property constructively. Therefore, there
is a unique canonical setoid isomorphism to every other fral.  Similarly,
by including the evaluation equations (\textsc{eval}), we obtain a
setoid frex together with a setoid isomorphism to every other frex.
These
isomorphisms let us extract simplification proofs generically out of
user-defined simplifiers. We discuss the decision to use setoids
in \secref{conclusion} where we refer to the concrete benefits setoids offer.
